\begin{lemma}
Let $\gamma_0,\gamma_1\colon[0,1]\to\mathbb{C}$ be continuous and nowhere zero,
and suppose they are homotopic relative to $\{0,1\}$ through
$\mathbb{C}-\{0\}$: that is, writing
$P=\{z\in\mathbb{C}:z\neq0\}$ with its subspace topology, there are continuous
maps $g_0,g_1\colon[0,1]\to P$ with $g_j(t)=\gamma_j(t)$ for all $t$, and a
homotopy from $g_0$ to $g_1$ inside $P$ which is stationary at $t=0$ and at
$t=1$.  Then
\[
  \operatorname{ind}(\gamma_0)=\operatorname{ind}(\gamma_1),
\]
with $\operatorname{ind}$ the index about $0$ of \noderef{WindingNumber}.
\end{lemma}
\begin{proof}
Write $\mathbb{T}=\mathbb{R}/2\pi\mathbb{Z}$ and let
$q\colon\mathbb{R}\to\mathbb{T}$ be the quotient map, a covering map.  Recall
from \noderef{WindingNumber} that for a continuous nowhere-zero
$\gamma\colon[0,1]\to\mathbb{C}$ the index is
\begin{equation}\tag{1}
  \operatorname{ind}(\gamma)=\frac{\Theta_\gamma(1)-\arg\gamma(0)}{2\pi},
\end{equation}
where $\Theta_\gamma$ is the lift along $q$ of the continuous map
$\alpha_\gamma\colon[0,1]\to\mathbb{T}$, $\alpha_\gamma(t)=q(\arg\gamma(t))$,
normalised by $\Theta_\gamma(0)=\arg\gamma(0)$.

\emph{Step 1: a continuous map $P\to\mathbb{T}$.}  Define
\[
  u\colon P\to\mathbb{T},\qquad u(z)=q(\arg z).
\]
This is continuous: the reduction $z\mapsto q(\arg z)$ of the principal
argument is continuous at every nonzero complex number, and every point of $P$
is such a point.  By construction, for $j\in\{0,1\}$ and every $t$,
\begin{equation}\tag{2}
  \bigl(u\circ g_j\bigr)(t)=q\bigl(\arg g_j(t)\bigr)=q\bigl(\arg\gamma_j(t)\bigr)
  =\alpha_{\gamma_j}(t),
\end{equation}
using $g_j(t)=\gamma_j(t)$.

\emph{Step 2: the endpoints agree.}  A homotopy relative to $\{0,1\}$ is
stationary at $0$ and at $1$, so its two ends agree there:
$g_0(0)=g_1(0)$ and $g_0(1)=g_1(1)$.  Hence
\begin{equation}\tag{3}
  \gamma_0(0)=\gamma_1(0),\qquad \gamma_0(1)=\gamma_1(1),
\end{equation}
and in particular $\arg\gamma_0(0)=\arg\gamma_1(0)$.  Call this common value
$e$, so that both lifts $\Theta_{\gamma_0}$ and $\Theta_{\gamma_1}$ are
normalised by the same initial value $e$.

\emph{Step 3: transport the homotopy.}  Post-composition with a continuous map
carries a homotopy relative to a set to a homotopy relative to the same set:
applying $u$ to the given homotopy from $g_0$ to $g_1$ produces a homotopy from
$u\circ g_0$ to $u\circ g_1$ relative to $\{0,1\}$, inside $\mathbb{T}$.  By
(2) this says
\[
  \alpha_{\gamma_0}\ \text{ and }\ \alpha_{\gamma_1}\ \text{ are homotopic relative to }\{0,1\}.
\]

\emph{Step 4: invariance of the lifted endpoint.}  For a covering map, two
paths in the base that are homotopic relative to $\{0,1\}$ have lifts with the
same terminal value, once the lifts are started at the same point: this is the
homotopy invariance of the endpoint of the lift.  Applying it to $q$, to the
two paths $\alpha_{\gamma_0}$ and $\alpha_{\gamma_1}$ of Step 3, and to the
common initial value $e$ of Step 2, we get
\[
  \Theta_{\gamma_0}(1)=\Theta_{\gamma_1}(1).
\]

\emph{Step 5: conclusion.}  By (1), (3) and Step 4,
\[
  \operatorname{ind}(\gamma_0)=\frac{\Theta_{\gamma_0}(1)-\arg\gamma_0(0)}{2\pi}
  =\frac{\Theta_{\gamma_1}(1)-\arg\gamma_1(0)}{2\pi}=\operatorname{ind}(\gamma_1).
\]
\end{proof}
